\section{Empirical Characterization of Constructive and Destructive Regimes}
\label{sec:gd-experiments}

We test how constructive and destructive regimes emerge during training,
how the gradient-based curvature estimate relates to their onset, and
how repeated reversal shapes accumulated parameter change. Full-batch
GD isolates these dynamics from sampling noise. We then vary batch size
to examine how stochasticity changes the observed alignment.

The comparison uses a tanh MLP on 5,000 CIFAR-10 images, ResNet-32 on all
50,000 CIFAR-10 training images, and a word-level Transformer on WikiText-2.
Each minimizes its fixed training objective for 3,000 updates.
The constant learning rates are $\{0.004,0.02,0.04\}$ for the MLP,
$\{0.025,0.05,0.1\}$ for ResNet, and $\{0.05,0.1,0.2\}$ for the Transformer.
Within each model, all rates share one frozen initialization and data
preparation. Prior training calibrations selected contrasting regimes;
test loss did not guide selection. Plain GD uses no momentum, weight
decay, clipping, or schedule. These three cases test recurrence across
objectives and architectures, conditional on their matched initializations.
Appendix~\ref{appendix:gd-raw-alignment} specifies the objectives and
measurement conventions.

\subsection{Curvature and Oscillatory Convergence}

Figure~\ref{fig:gd-loss} shows a clear separation between smooth and
oscillatory convergence. The lowest MLP and Transformer rates produce
smooth loss curves and $\eta\widehat\kappa_t$ near zero. At higher rates,
the loss develops fluctuations while its longer-term trend continues
downward; the estimated scaled curvature approaches two. ResNet enters
this regime at every tested rate, with earlier transitions at larger
learning rates.

The retained Hessian measurements place these signatures in their
spectral context. As sharpness grows, the trajectories develop the
oscillations associated with EoS, while the inexpensive gradient
diagnostic identifies the curvature response along the direction
actually followed. In the oscillatory regime, $\eta\widehat\kappa_t$
remains close to the local loss-change boundary even when
$\eta\lambda_{\max}$ is larger. The distinction is expected: directional
curvature depends on the gradient's orientation, whereas the leading
eigenvalue describes the sharpest direction. ResNet and Transformer
spectra additionally use fixed subsets of their training objectives.
The comparison therefore supports monitoring the realized curvature
response throughout training, with no additional Hessian computation.

\subsection{Destructive Regimes and Displacement Retention}

The loss oscillations coincide with a transition to destructive
alignment in Fig.~\ref{fig:gd-alignment-displacement}. The lowest MLP and
Transformer rates retain positive cosine, while the other seven
conditions develop nearly opposing consecutive gradients. In the MLP
and ResNet, this transition also accompanies growing gradient norms:
larger steps increasingly reverse preceding parameter changes. The
joint behavior directly exposes the magnitude and directional factors
in Eq.~\eqref{eq:loss-rise-alignment}.

Normalized displacement then declines strongly. The larger steps
accumulate path length faster than net parameter change, so $R_t$
falls even as training continues to reduce the loss. At the highest
rates, the retained fraction after 3,000 updates is approximately
2.6\% for the MLP, 3.7\% for ResNet, and 5.4\% for the Transformer.
Nevertheless, each highest-rate condition reaches its model's lowest
final training loss. Destructive regimes therefore sustain convergence
despite extensive local reversal. The risk of an upward loss excursion
depends on how strongly gradient growth amplifies that reversal.

\subsection{Oscillation across Temporal Lags}

We next ask whether destructive alignment reflects an isolated pair
of updates or a recurring temporal structure. Figure~\ref{fig:gd-backward}
fixes the gradient after 2,256 updates and compares it with the preceding
256 gradients, using the highest rate for each model. The first lags
alternate between strongly negative odd-lag and strongly positive
even-lag alignment. This is direct evidence of repeated back-and-forth
motion: a gradient can oppose its immediate predecessor while reinforcing
the direction from two steps earlier.

The early oscillation envelope contracts, rapidly for the MLP and ResNet
and more slowly for the Transformer. Across the full window, alignment
can strengthen again, so the envelope is not monotone. These recurrences
explain how local destructive interference coexists with reinforcement
over longer lags and why net progress survives repeated reversals.
